\section{Permutation-Based Physical Symmetry}

The linear-combination twist $\widetilde{A}_g^i = \sum_j U(g)_{ij}\, A^j$
is the standard form used in the virtual representation theorem.
The following alternative formulation replaces the matrix action by a
permutation of the physical index; it is useful for symmetries that act by
reordering basis states rather than by a general linear map. % The statements of this section are auxiliary constructions not
% directly sourced from~\cite{PerezGarcia2008StringOrder}.

\begin{definition}[Permutation symmetry datum]\label{def:on_site_symmetry_perm}
    \lean{MPSTensor.OnSiteSymmetry}
    \leanok
    An \emph{on-site symmetry datum} for a group $G$ on a physical
    space of dimension $d$ consists of a matrix-valued map
    $u : G \to \MN{d}$ and a physical-index action
    $\sigma : G \to \mathrm{End}(\{0,\ldots,d{-}1\})$.
\end{definition}

\begin{definition}[Permutation-twisted tensor]\label{def:perm_twisted_tensor}
    \lean{MPSTensor.TwistedTensor}
    \leanok
    \uses{def:on_site_symmetry_perm, def:mps_tensor}
    Given an on-site symmetry datum $(u, \sigma)$, the
    \emph{permutation-twisted tensor} at group element $g$ is
    \begin{align}
        (A^{\sigma_g})^i &:= A^{\sigma(g)(i)}.
        \label{eq:symmetry_perm_twist}
    \end{align}
    In tensor-network notation the local tensor is unchanged and only the
    physical leg is relabelled:
    % Formula: (A^{\sigma_g})^i = A^{\sigma(g)(i)}.
    % Source verdict: the adjacent defining identity is the source; there is
    % no dedicated source figure.
    % Ink-to-index: i and \sigma(g)(i) label the free upper physical stub of
    % the same one-site tensor A; the relabelling arrow carries no index.
    % Contractions: none.
    % Paired boundary: boundary|virtual-west=1|virtual-east=1|
    % physical-up=1|physical-down=0.
    \begin{align}
        \begin{tenkz}[physical=up]
            \tn[up=$i$]{A}
        \end{tenkz}
        &\qquad \xrightarrow{\sigma_g} \qquad
        \begin{tenkz}[physical=up]
            \tn[up=$\sigma(g)(i)$]{A}
        \end{tenkz}.
        \notag
    \end{align}
\end{definition}

\begin{lemma}[Identity permutation twist]\label{lem:perm_twisted_one}
    \lean{MPSTensor.TwistedTensor_one}
    \leanok
    \uses{def:perm_twisted_tensor}
    If $\sigma(1) = \id$, then $(A^{\sigma_1})^i = A^i$ for all~$i$.
\end{lemma}

\begin{proof}\leanok\uses{def:perm_twisted_tensor}
    The claim follows from $\sigma(1)(i) = i$ for all~$i$.
\end{proof}

\noindent\emph{Composition order.}
The linear-combination twist (Lemma~\ref{lem:twisted_tensor_mul}) composes as
$\widetilde{(\widetilde{A}_h)}_g = \widetilde{A}_{gh}$, applying $h$ first
then~$g$. The permutation twist below uses the same left-action convention
$\sigma(gh) = \sigma(g) \circ \sigma(h)$, so
$(A^{\sigma_g})^{\sigma_h} = A^{\sigma_{gh}}$; the inner permutation
$\sigma(h)$ acts first on the index.

\begin{lemma}[Composition of permutation twists]\label{lem:perm_twisted_mul}
    \lean{MPSTensor.TwistedTensor_mul}
    \leanok
    \uses{def:perm_twisted_tensor}
    If $\sigma(gh) = \sigma(g) \circ \sigma(h)$ for all $g, h \in G$,
    then
    \begin{align}
        A^{\sigma_{gh}} &= (A^{\sigma_g})^{\sigma_h}.
        \label{eq:symmetry_perm_composition}
    \end{align}
\end{lemma}

\begin{proof}\leanok\uses{def:perm_twisted_tensor}
    Expanding both sides gives
    $(A^{\sigma_{gh}})^i
    = A^{\sigma(g)(\sigma(h)(i))}
    = ((A^{\sigma_g})^{\sigma_h})^i$.
    Diagrammatically this is just two consecutive relabellings of the same
    local tensor:
    % Formula: A^i \xrightarrow{\sigma_h} A^{\sigma(h)(i)}
    %   \xrightarrow{\sigma_g} A^{\sigma(g)(\sigma(h)(i))}.
    % Source verdict: the adjacent componentwise equality is the source;
    % there is no dedicated source figure.
    % Ink-to-index: the three upper stubs carry, from left to right, the free
    % physical indices i, \sigma(h)(i), and \sigma(g)(\sigma(h)(i)).  The two
    % arrows record relabelling by \sigma_h and then \sigma_g; they are not
    % tensor-network contractions.
    % Contractions: none in any panel.
    % Common boundary: boundary|virtual-west=1|virtual-east=1|
    % physical-up=1|physical-down=0.
    \begin{align}
        \tnpic[physical=up]{\tn[up=$i$]{A}}
        &\qquad \xrightarrow{\sigma_h} \qquad
        \tnpic[physical=up]{\tn[up=$\sigma(h)(i)$]{A}}
        \qquad \xrightarrow{\sigma_g} \qquad
        \tnpic[physical=up]{\tn[up=$\sigma(g)(\sigma(h)(i))$]{A}}.
        \notag
    \end{align}
\end{proof}

\section{String Order and Local Symmetry Equivalence}

The string order parameter~\cite{PerezGarcia2008StringOrder} detects
whether an on-site unitary $u$ acts as a local symmetry on the state
generated by an injective MPS tensor.

\begin{definition}[Twisted transfer map]\label{def:twisted_transfer_map}
    \lean{MPSTensor.twistedTransferMap}
    \leanok
    \uses{def:mps_tensor}
    Given an MPS tensor $A$ and a physical-index matrix $u$, the
    \emph{twisted transfer map} is
    \begin{align}
        \E_u(X)
        &:= \sum_{n,n'} \bra{n'} u \ket{n}
        A_n\, X\, A_{n'}^\dagger.
        \label{eq:symmetry_twisted_transfer}
    \end{align}
    Diagrammatically one inserts the physical operator $u$ on the doubled
    local transfer cell:
    % Applied formula: \mathcal E_u(X)
    % = \sum_{n,n'}\langle n'|u|n\rangle A_n X A_{n'}^\dagger.
    % Equivalent unapplied doubled formula: \mathcal E_u
    % = \sum_{n,n'}\langle n'|u|n\rangle
    % A_n\otimes\overline{A_{n'}}.
    % Source verdict: Papers/0802.0447/StringOrder-v10.tex, equation EU,
    % lines 164--174, gives this exact algebraic identity but no dedicated
    % figure.  The lower \tn* supplies entrywise conjugation; reading order
    % supplies the transpose that yields A_{n'}^\dagger in the applied map.
    % Pairing/index map: the top-to-middle pairleg contracts n and the
    % middle-to-bottom pairleg contracts n'.  These are the only physical
    % pairings.  The middle u row has no virtual boundary, while both
    % virtual legs of A_n and both virtual legs of \overline{A_{n'}} remain
    % free because the map has not consumed X.
    % boundary|virtual-west=2|virtual-east=2|physical-up=0|physical-down=0.
    \begin{align}
        &\begin{tenkz}[rows={ket,op:none,bra}]
            \tn{A_n} \\
            \tn[role=operator]{u} \\
            \tn*{A_{n'}}
        \end{tenkz}.
        \notag
    \end{align}
    The upper and lower black nodes denote $A_n$ and $A_{n'}^\dagger$, and the
    red node labelled $u$ couples the physical legs.
\end{definition}

\begin{definition}[String order parameter]\label{def:string_order_param}
    \lean{MPSTensor.stringOrderParam}
    \leanok
    \uses{def:twisted_transfer_map}
    The \emph{string order parameter} for twist $u$ and stationary
    state $\Lambda$ at length $L$ is
    \begin{align}
        R_L(u) &:= \tr\!(\Lambda \cdot \E_u^L(\Id)).
        \label{eq:symmetry_string_parameter}
    \end{align}
    In tensor-network notation this is the length-$L$ doubled chain with a
    copy of $u$ inserted on each physical rung:
    % Formula/source: Papers/0802.0447/StringOrder-v10.tex,
    % lines 380--385 explicitly define R_L; equations EU and SOPMP,
    % lines 164--181, give the constituent transfer/MPS expressions. The
    % displayed transfer formula follows after canonical fixed-point/unital reduction.
    % Ink-to-index map: each ket physical index contracts into u and each
    % output of u contracts into the matching bra index; no physical leg is free.
    % Contraction set: all on-site ket--u--bra rungs, all horizontal
    % ket/bra bonds, and both boundary cups are summed; Lambda labels the west cup.
    % Type verdict: a scalar after every physical and virtual index is contracted.
    % Boundary signature: virtual-west=0 | virtual-east=0 |
    % physical-up=0 | physical-down=0.
    \begin{align}
        &\begin{tenkz}[rows={ket,op:none,bra},
                      west={cup=$\Lambda$}, east=cup]
            \tn{A} & \tn{A} & \tndots & \tn{A} \\
            \tn[role=operator]{u} & \tn[role=operator]{u} & \tndots &
                \tn[role=operator]{u} \\
            \tn*{A}\tnspan[brace below]{4}{$L\text{ sites}$} &
                \tn*{A} & \tndots & \tn*{A}
        \end{tenkz}.
        \notag
    \end{align}
    The boundary matrix $\Lambda$ is contracted into the virtual boundary, the
    red nodes labelled $u$ sit on the physical rungs, and the right boundary
    is traced.
\end{definition}

\begin{definition}[Physical string correlator]\label{def:physical_string_order_param}
    \lean{MPSTensor.physicalStringOrderParam}
    \leanok
    \uses{def:twisted_transfer_map}
    For physical operators $x,y \in \MN{d}$, a twist $u \in \MN{d}$, and
    boundary state $\Lambda$, the transfer-form string correlator of
    \cite[display~SOPMP, lines~176--181]{PerezGarcia2008StringOrder} is
    \begin{align}
        S_N(x,y,u)
        &:= \tr\!(\Lambda\,\E_x\,\E_u^N(\E_y(\Id))).
        \label{eq:symmetry_physical_string}
    \end{align}
    This is the physical-endpoint expression corresponding to
    $\langle \Psi|\,x \otimes u^{\otimes N} \otimes y\,|\Psi\rangle$.
\end{definition}

\begin{definition}[Physical string order]\label{def:has_physical_string_order}
    \lean{MPSTensor.HasPhysicalStringOrderWith}
    \lean{MPSTensor.HasPhysicalStringOrder}
    \leanok
    \uses{def:physical_string_order_param}
    A tensor has physical string order, in the sense of
    \cite[display~SOP, lines~112--122]{PerezGarcia2008StringOrder}, if
    there exist a unitary $u \ne \Id$ and physical endpoint operators $x,y$
    for which the absolute value of the correlator has a positive limit:
    for some $s>0$, $|S_N(x,y,u)| \longrightarrow s$.
    The comparison with virtual boundary matrices is a separate
    endpoint-realizability statement.
\end{definition}

\begin{definition}[Boundary string order parameter]%
    \label{def:string_order_boundary_param}
    \lean{MPSTensor.stringOrderBoundaryParam}
    \leanok
    \uses{def:twisted_transfer_map, def:string_order_param}
    For a boundary state $\Lambda$ and boundary matrices $X,Y \in \MN{D}$, the
    \emph{boundary string order parameter} is
    \begin{align}
        R_L(u;X,Y)
        &:= \tr\!(\Lambda X \E_u^L(Y)).
        \label{eq:symmetry_boundary_string}
    \end{align}
    The ordinary string order parameter is the special case $X=Y=\Id$.
    The matrices $X,Y$ act on the virtual level. In the string correlator
    of~\cite{PerezGarcia2008StringOrder} the string of twists is instead
    flanked by local operators $x,y$ acting on physical sites; expanding
    such operators in the transfer picture produces particular virtual
    matrices, and \cite{PerezGarcia2008StringOrder} argues that for
    injective tensors, physical operators on sufficiently many sites
    produce every virtual matrix. The results below are stated directly
    in this virtual-boundary form.
    % Deviation recorded in docs/paper-gaps/pgwsvc08_string_order_virtual_boundary.tex.
\end{definition}

\begin{definition}[Local symmetry]\label{def:local_symmetry}
    \lean{MPSTensor.IsLocalSymmetry}
    \leanok
    \uses{def:mps_tensor}
    An MPS tensor $A$ has \emph{local symmetry} under a unitary $u$ with
    respect to a boundary state $\Lambda$ if there exist a unitary
    matrix $V$ and a phase $\mu$ with $|\mu| = 1$ such that
    $V^\dagger \Lambda V = \Lambda$ and, for every physical index~$i$,
    \begin{align}
        \sum_j u_{ij}\, A^j
        &= \mu\, V\, A^i\, V^\dagger.
        \label{eq:symmetry_local_covariance}
    \end{align}
    The unitary $V$ intertwines the physical action of~$u$ with a virtual
    conjugation, and the boundary state $\Lambda$ is invariant under~$V$.
\end{definition}

\begin{definition}[String order]\label{def:has_string_order}
    \lean{MPSTensor.HasStringOrder}
    \leanok
    \uses{def:string_order_boundary_param}
    String order \emph{exists} for $u$ with respect to a boundary
    state $\Lambda$ if there exist boundary matrices $X, Y$ and a
    positive constant $c > 0$ such that
    $c \le \|R_L(u;X,Y)\|$ for all $L$, where
    $R_L(u;X,Y)$ is the boundary-refined string-order parameter.
    Thus some boundary-refined string-order sequence does not decay to zero.

    String order is here an existence statement over the virtual matrices
    $X,Y$, for a fixed twist $u$; no claim is made for any particular
    boundary choice. Definition~\ref{def:has_physical_string_order} records
    the physical-endpoint source definition, which quantifies existentially
    over $u,x,y$ and requires a positive limit rather than a lower bound
    uniform in the length; that paper also notes that the correlator for a
    particular endpoint choice can vanish even when the symmetry is present.
    % Deviation recorded in docs/paper-gaps/pgwsvc08_string_order_virtual_boundary.tex.
\end{definition}

\begin{lemma}[String order prevents decay to zero]
    \label{lem:not_tendsto_zero_of_hasStringOrder}
    \lean{MPSTensor.not_tendsto_zero_of_hasStringOrder}
    \leanok
    \uses{def:has_string_order}
    If $A$ has string order for $u$ with respect to a boundary state
    $\Lambda$, then for some boundary matrices $X,Y$ the sequence
    $R_L(u;X,Y)$ does not tend to $0$ as $L \to \infty$.
\end{lemma}

\begin{proof}
    \leanok
    \uses{def:has_string_order}
    A uniform lower bound $c \le \|R_L(u;X,Y)\|$ is incompatible with
    convergence to $0$.
\end{proof}

\begin{definition}[Local physical intertwining]\label{def:condC1}
    \lean{MPSTensor.CondC1}
    \leanok
    \uses{def:mps_tensor}
    The local physical intertwining relation is
    $\sum_j u_{ij}\, A^j = V\, A^i\, V^\dagger$.
    Locally, the physical action of $u$ can therefore be pushed to the virtual
    legs:
    \begin{tenkzequation}
        % Formula/source: Papers/0802.0447/StringOrder-v10.tex,
        % Condition C1, lines 328--334, states
        % sum_j u_{ij} A^j = V A^i V^\dagger.
        % Ink-to-index map: i is free above u; j is the sole physical
        % contraction between u and A; west/east virtual indices remain free.
        % Contraction set: one physical u--A bond on the left and two
        % virtual gauge bonds on the right.
        % Type verdict: both sides are rank-three local tensors.
        % Boundary signature: virtual-west=1 | virtual-east=1 |
        % physical-up=1 | physical-down=0 on both sides.
        \begin{tenkz}[rows={op:none,ket}]
            \tn[role=operator,up=$i$]{u} \\
            \tn{A}
        \end{tenkz}
        =
        % Formula/source: Papers/0802.0447/StringOrder-v10.tex,
        % Condition C1, lines 328--334, states
        % sum_j u_{ij} A^j = V A^i V^\dagger.
        % Ink-to-index map: i is free above u; j is the sole physical
        % contraction between u and A; west/east virtual indices remain free.
        % Contraction set: one physical u--A bond on the left and two
        % virtual gauge bonds on the right.
        % Type verdict: both sides are rank-three local tensors.
        % Boundary signature: virtual-west=1 | virtual-east=1 |
        % physical-up=1 | physical-down=0 on both sides.
        \begin{tenkz}[physical=up]
            \tnX{V} & \tn[up=$i$]{A} & \tnX{V^\dagger}
        \end{tenkz}.
    \end{tenkzequation}
\end{definition}

\begin{definition}[Transfer-map covariance]\label{def:condC2}
    \lean{MPSTensor.CondC2}
    \leanok
    \uses{def:transfer_map}
    Transfer-map covariance is the identity that the transfer map commutes with
    virtual conjugation:
    \begin{align}
        \E(V X V^\dagger)
        &= V\, \E(X)\, V^\dagger.
        \label{eq:symmetry_transfer_covariance}
    \end{align}
    In doubled-layer notation, the virtual operators slide through the local
    transfer cell:
    \begin{tenkzequation}
        % Formula/source: Papers/0802.0447/StringOrder-v10.tex,
        % Condition C2, lines 335--340, states E(VXV^\dagger)=V E(X)V^\dagger.
        % Ink-to-index map: the ket row carries V and A, the bra row carries
        % their complex conjugates; each column of A and bar(A) shares one index.
        % Contraction set: one physical ket--bra pairing per transfer cell and
        % the internal horizontal bonds on each row; all four outer wires are free.
        % Type verdict: equality of endomorphisms on the doubled virtual space.
        % Boundary signature: virtual-west=2 | virtual-east=2 |
        % physical-up=0 | physical-down=0 on both sides.
        \begin{tenkz}[sandwich]
            \tnX{V} & \tn{A} \\
            \tnX{\overline V} & \tn*{A}
        \end{tenkz}
        =
        % Formula/source: Papers/0802.0447/StringOrder-v10.tex,
        % Condition C2, lines 335--340, states E(VXV^\dagger)=V E(X)V^\dagger.
        % Ink-to-index map: the ket row carries V and A, the bra row carries
        % their complex conjugates; each column of A and bar(A) shares one index.
        % Contraction set: one physical ket--bra pairing per transfer cell and
        % the internal horizontal bonds on each row; all four outer wires are free.
        % Type verdict: equality of endomorphisms on the doubled virtual space.
        % Boundary signature: virtual-west=2 | virtual-east=2 |
        % physical-up=0 | physical-down=0 on both sides.
        \begin{tenkz}[sandwich]
            \tn{A} & \tnX{V} \\
            \tn*{A} & \tnX{\overline V}
        \end{tenkz}.
    \end{tenkzequation}
\end{definition}

\begin{definition}[Commutation in the doubled transfer picture]\label{def:condC3}
    \lean{MPSTensor.CondC3}
    \leanok
    \uses{def:twisted_transfer_map}
    In the transfer-map language, the commutation of the doubled transfer matrix
    $E = \sum_i A_i \otimes \bar A_i$ with $V \otimes \bar V$ is
    \begin{align}
        V\, \E_1(X)\, V^\dagger
        &= \E_1(V X V^\dagger).
        \label{eq:symmetry_doubled_commutation}
    \end{align}
\end{definition}

\begin{theorem}[Transfer-map covariance and doubled commutation are equivalent]%
    \label{thm:condC2_iff_condC3}
    \lean{MPSTensor.condC2_iff_condC3}
    \leanok
    \uses{def:condC2, def:condC3}
    Transfer-map covariance and commutation in the doubled transfer
    picture are equivalent for any matrix $V$.
\end{theorem}

\begin{proof}
    \leanok
    Both sides express the same local picture with opposite orientation:
    \begin{tenkzequation}
        % Formula/source: Papers/0802.0447/StringOrder-v10.tex,
        % Conditions C2--C3, lines 335--347, identify covariance with
        % commutation by V\otimes\overline V in the doubled picture.
        % Ink-to-index map: the upper/lower rows carry the ket/bra virtual
        % indices; the A--bar(A) pair contracts the physical Kraus index.
        % Contraction set: the transfer-cell physical pairing and each internal
        % horizontal row bond are summed; the four outer virtual indices are free.
        % Type verdict: equality of two doubled-space endomorphisms.
        % Boundary signature: virtual-west=2 | virtual-east=2 |
        % physical-up=0 | physical-down=0 on both sides.
        \begin{tenkz}[sandwich]
            \tnX{V} & \tn{A} \\
            \tnX{\overline V} & \tn*{A}
        \end{tenkz}
        =
        % Formula/source: Papers/0802.0447/StringOrder-v10.tex,
        % Conditions C2--C3, lines 335--347, identify covariance with
        % commutation by V\otimes\overline V in the doubled picture.
        % Ink-to-index map: the upper/lower rows carry the ket/bra virtual
        % indices; the A--bar(A) pair contracts the physical Kraus index.
        % Contraction set: the transfer-cell physical pairing and each internal
        % horizontal row bond are summed; the four outer virtual indices are free.
        % Type verdict: equality of two doubled-space endomorphisms.
        % Boundary signature: virtual-west=2 | virtual-east=2 |
        % physical-up=0 | physical-down=0 on both sides.
        \begin{tenkz}[sandwich]
            \tn{A} & \tnX{V} \\
            \tn*{A} & \tnX{\overline V}
        \end{tenkz}.
    \end{tenkzequation}
    The two formulas are the same identity read from opposite sides.
\end{proof}

\begin{lemma}[Unitary Kraus mixing]\label{lem:unitary_kraus_mixing}
    \lean{MPSTensor.unitary_kraus_mixing}
    \leanok
    If $u$ is unitary, then replacing Kraus operators $\{A_i\}$ by
    their unitary mixtures $\{\sum_j u_{ij} A_j\}$ preserves the
    channel:
    \begin{align}
        \sum_i \left(\sum_j u_{ij} A_j\right)\, Y
        \left(\sum_j u_{ij} A_j\right)^\dagger
        &=
        \sum_i A_i\, Y\, A_i^\dagger.
        \label{eq:symmetry_unitary_kraus_mixing}
    \end{align}
\end{lemma}

\begin{proof}
    \leanok
    \uses{thm:kraus_same_map_of_unitary_combination}
    This is exactly the usual invariance of a Kraus decomposition under a
    unitary change of Kraus operators; see
    Theorem~\ref{thm:kraus_same_map_of_unitary_combination}.
\end{proof}

\begin{theorem}[Local physical intertwining implies transfer-map covariance]%
    \label{thm:condC1_imp_condC2}
    \lean{MPSTensor.condC1_imp_condC2}
    \leanok
    \uses{def:condC1, def:condC2, lem:unitary_kraus_mixing}
    If $V$ and $u$ are unitary and the local physical intertwining
    relation holds, then the transfer map is covariant.
\end{theorem}

\begin{proof}
    \leanok
    \uses{lem:unitary_kraus_mixing}
    Starting from
    \begin{tenkzequation}
        % Formula/source: Papers/0802.0447/StringOrder-v10.tex,
        % Condition C1, lines 328--334, is the displayed intertwining relation.
        % Ink-to-index map: i stays free, j contracts vertically between u and A,
        % and the west/east virtual indices stay free on both sides.
        % Contraction set: one physical bond on the left; two virtual gauge
        % bonds on the right.
        % Type verdict: equality of rank-three local tensors.
        % Boundary signature: virtual-west=1 | virtual-east=1 |
        % physical-up=1 | physical-down=0 on both sides.
        \begin{tenkz}[rows={op:none,ket}]
            \tn[role=operator,up=$i$]{u} \\
            \tn{A}
        \end{tenkz}
        =
        % Formula/source: Papers/0802.0447/StringOrder-v10.tex,
        % Condition C1, lines 328--334, is the displayed intertwining relation.
        % Ink-to-index map: i stays free, j contracts vertically between u and A,
        % and the west/east virtual indices stay free on both sides.
        % Contraction set: one physical bond on the left; two virtual gauge
        % bonds on the right.
        % Type verdict: equality of rank-three local tensors.
        % Boundary signature: virtual-west=1 | virtual-east=1 |
        % physical-up=1 | physical-down=0 on both sides.
        \begin{tenkz}[physical=up]
            \tnX{V} & \tn[up=$i$]{A} & \tnX{V^\dagger}
        \end{tenkz},
    \end{tenkzequation}
    one inserts $V^\dagger V = \Id$ so that each local transfer cell is
    written in terms of conjugated Kraus operators, then replaces those
    operators using the intertwining relation. After that, the physical
    $u$-boxes disappear by the unitary Kraus-mixing identity.
\end{proof}

\begin{theorem}[Transfer-map covariance implies local physical intertwining]%
    \label{thm:condC2_imp_condC1_of_injective}
    \lean{MPSTensor.condC2_imp_condC1_of_injective}
    \leanok
    \uses{def:condC2, def:condC1, def:injective}
    For an injective MPS, if $V$ is unitary and the transfer map is
    covariant, then there exists a unitary $u$ satisfying the local
    physical intertwining relation.
\end{theorem}

\begin{proof}
    \leanok
    \uses{thm:condC2_iff_condC3, thm:condC1_imp_condC2}
    Set $B^i := V\, A^i\, V^\dagger$. The transfer-map covariance relation
    \begin{tenkzequation}
        % Formula/source: Papers/0802.0447/StringOrder-v10.tex,
        % Condition C2, lines 335--340, is E\circ Ad_V=Ad_V\circ E.
        % Ink-to-index map: V and bar(V) act on the ket/bra virtual rows;
        % A and bar(A) share the summed physical Kraus index.
        % Contraction set: one transfer-cell physical pairing and both internal
        % horizontal row bonds; the four boundary wires remain free.
        % Type verdict: equality of doubled-space endomorphisms.
        % Boundary signature: virtual-west=2 | virtual-east=2 |
        % physical-up=0 | physical-down=0 on both sides.
        \begin{tenkz}[sandwich]
            \tnX{V} & \tn{A} \\
            \tnX{\overline V} & \tn*{A}
        \end{tenkz}
        =
        % Formula/source: Papers/0802.0447/StringOrder-v10.tex,
        % Condition C2, lines 335--340, is E\circ Ad_V=Ad_V\circ E.
        % Ink-to-index map: V and bar(V) act on the ket/bra virtual rows;
        % A and bar(A) share the summed physical Kraus index.
        % Contraction set: one transfer-cell physical pairing and both internal
        % horizontal row bonds; the four boundary wires remain free.
        % Type verdict: equality of doubled-space endomorphisms.
        % Boundary signature: virtual-west=2 | virtual-east=2 |
        % physical-up=0 | physical-down=0 on both sides.
        \begin{tenkz}[sandwich]
            \tn{A} & \tnX{V} \\
            \tn*{A} & \tnX{\overline V}
        \end{tenkz}
    \end{tenkzequation}
    implies that the two Kraus families $B$ and $A$ define the same
    transfer map. Kraus freedom for equal channels converts
    $\E_B = \E_A$ into a unitary Kraus mixing matrix $u$ with
    $B^i = \sum_j u_{ij} A^j$.
\end{proof}

\begin{definition}[Twisted companion tensor]\label{def:twisted_mixed_companion}
    \lean{MPSTensor.twistedMixedCompanion}
    \leanok
    \uses{def:mps_tensor}
    For an MPS tensor $A$ and a physical-index matrix $u$, define the
    companion tensor $B_u$ by
    \begin{align}
        B_u^n
        &:= \sum_{n'=0}^{d-1} \overline{u_{n'n}}\, A^{n'}.
        \label{eq:symmetry_companion_tensor}
    \end{align}
\end{definition}

\begin{lemma}[Twisted transfer as a mixed transfer operator]%
    \label{lem:twisted_transfer_eq_mixed_transfer}
    \lean{MPSTensor.twistedTransferMap_eq_mixedTransfer}
    \leanok
    \uses{def:twisted_transfer_map, def:twisted_mixed_companion, def:mixed_transfer}
    The twisted transfer map of $A$ is the mixed transfer operator of the pair
    $(A,B_u)$:
    \begin{align}
        \E_u &= F_{A,B_u}.
        \label{eq:symmetry_twisted_mixed_transfer}
    \end{align}
\end{lemma}

\begin{proof}
    \leanok
    Expand both definitions and rearrange the double sum.
\end{proof}

\begin{theorem}[Twisted-transfer eigenvalues have modulus at most one]%
    \label{thm:twisted_transfer_spectral_radius_le_one}
    \lean{MPSTensor.twistedTransfer_spectralRadius_le_one}
    \leanok
    \uses{def:twisted_transfer_map, def:injective}
    Let $A$ be an injective MPS tensor with $\E_A(\Id)=\Id$, and let $u$
    satisfy $u u^\dagger = \Id$.
    If $\E_u(X) = \lambda X$ for some nonzero $X \in \MN{D}$, then
    $|\lambda| \le 1$.
\end{theorem}

\begin{proof}
    \leanok
    \uses{lem:unitary_kraus_mixing, lem:twisted_transfer_eq_mixed_transfer,
        thm:left_canonical_gauge, thm:eigenvalue_bound}
    By Lemma~\ref{lem:twisted_transfer_eq_mixed_transfer}, the twisted transfer
    map is the mixed transfer operator $F_{A,B_u}$, as in
    (\ref{eq:symmetry_twisted_mixed_transfer}).
    Because $B_u$ is obtained from $A$ by unitary mixing of the Kraus operators,
    the transfer maps of $A$ and $B_u$ coincide, so the two adjoint channels
    share a common positive definite fixed point.
    Gauging both families by that fixed point puts them in a common
    trace-preserving gauge.
    Theorem~\ref{thm:eigenvalue_bound} then applies to the mixed transfer
    operator in that gauge, and similarity preserves the eigenvalue $\lambda$.
\end{proof}

\begin{theorem}[Modulus-one twisted eigenvalue implies gauge-phase equivalence]%
    \label{thm:twisted_transfer_modulus_one_gauge_phase}
    \lean{MPSTensor.twistedTransfer_modulus_one_implies_gaugePhase}
    \leanok
    \uses{def:twisted_transfer_map, def:twisted_mixed_companion,
        def:gauge_phase_equiv, def:injective}
    Let $A$ be an injective MPS tensor with $\E_A(\Id)=\Id$, and let $u$
    satisfy $u u^\dagger = \Id$.
    If
    \begin{align}
        \E_u(X) &= \lambda X
        \label{eq:symmetry_twisted_eigen_mod_one}
    \end{align}
    for some nonzero $X \in \MN{D}$ with $|\lambda|=1$, then $A$ is
    gauge-phase equivalent to the companion tensor $B_u$.
\end{theorem}

\begin{proof}
    \leanok
    \uses{lem:unitary_kraus_mixing, lem:twisted_transfer_eq_mixed_transfer,
        thm:left_canonical_gauge, thm:modulus_one_gauge_irr_tp}
    By Lemma~\ref{lem:twisted_transfer_eq_mixed_transfer}, the twisted transfer
    map is the mixed transfer operator $F_{A,B_u}$.
    Because $B_u$ is obtained from $A$ by unitary mixing of the Kraus operators,
    the transfer maps of $A$ and $B_u$ coincide, so the two adjoint channels
    share a common positive definite fixed point.
    Gauging both families by that fixed point puts them in a common
    trace-preserving gauge.
    In that gauge, (\ref{eq:symmetry_twisted_eigen_mod_one}) becomes a
    modulus-one peripheral eigenvalue for an irreducible mixed transfer
    operator, so Theorem~\ref{thm:modulus_one_gauge_irr_tp} yields
    gauge-phase equivalence.
    Undoing the common trace-preserving gauge gives the claim for $A$ and $B_u$.
\end{proof}

\begin{lemma}[Spectral radius $<1$ forces boundary-string decay]%
    \label{lem:boundary_string_order_tendsto_zero}
    \lean{MPSTensor.stringOrderBoundaryParam_tendsto_zero_of_spectralRadius_lt_one}
    \leanok
    \uses{def:string_order_boundary_param, def:twisted_transfer_map}
    If $\rho(\E_u) < 1$, then, as $L \to \infty$, for all boundary matrices
    $X,Y,\Lambda \in \MN{D}$ one has
    $R_L(u;X,Y) \longrightarrow 0$.
\end{lemma}

\begin{proof}
    \leanok
    In finite dimension, $\rho(\E_u) < 1$ implies $\E_u^L \to 0$ in operator
    norm. Composing with the fixed trace pairing
    $Z \mapsto \tr(\Lambda X Z)$ gives the stated scalar convergence.
\end{proof}

\begin{theorem}[String order gives gauge-phase equivalence with the companion tensor]%
    \label{thm:string_order_implies_twisted_gauge_phase}
    \lean{MPSTensor.gaugePhaseEquiv_twisted_of_hasStringOrder}
    \leanok
    \uses{def:has_string_order, def:twisted_mixed_companion,
        def:gauge_phase_equiv, def:injective}
    Let $A$ be injective, let $\Lambda \in \MN{D}$, and assume
    $u u^\dagger = \Id$ and $\E_A(\Id)=\Id$.
    If string order exists for $u$ with boundary parameter $\Lambda$, then $A$
    is gauge-phase equivalent to the companion tensor $B_u$.
\end{theorem}

\begin{proof}
    \leanok
    \uses{lem:boundary_string_order_tendsto_zero,
        thm:twisted_transfer_spectral_radius_le_one,
        thm:twisted_transfer_modulus_one_gauge_phase}
    If $\rho(\E_u) < 1$, Lemma~\ref{lem:boundary_string_order_tendsto_zero}
    forces every boundary string-order sequence to decay to zero, contradicting
    Definition~\ref{def:has_string_order}.
    Hence $\rho(\E_u) \ge 1$.
    Theorem~\ref{thm:twisted_transfer_spectral_radius_le_one} shows that every
    eigenvalue has modulus at most $1$, so some peripheral eigenvalue has
    modulus exactly $1$.
    Applying Theorem~\ref{thm:twisted_transfer_modulus_one_gauge_phase} then
    gives the gauge-phase equivalence.
\end{proof}

\begin{theorem}[Gauge-phase equivalence with the companion tensor yields a virtual unitary]%
    \label{thm:virtual_unitary_of_twisted_gauge_phase}
    \lean{MPSTensor.virtualUnitary_of_gaugePhaseEquiv_twisted}
    \leanok
    \uses{def:gauge_phase_equiv, def:twisted_mixed_companion, def:injective}
    Let $A$ be injective, assume $u u^\dagger = \Id$ and $\E_A(\Id)=\Id$.
    If $A$ is gauge-phase equivalent to the companion tensor $B_u$, then there
    exist a unitary $V$ and a phase $\mu$ with $|\mu|=1$ such that, for every
    physical index~$i$,
    \begin{align}
        \sum_j u_{ij} A^j
        &= \mu\, V A^i V^\dagger.
        \label{eq:symmetry_phased_covariance}
    \end{align}
\end{theorem}

\begin{proof}
    \leanok
    \uses{thm:psd_fp_unique}
    Start from an invertible intertwiner $X$ and a nonzero gauge scalar
    $\zeta$ between $A$ and $B_u$. The positive matrix
    $Q:=XX^\dagger$ satisfies $\E_A(Q)=|\zeta|^2Q$. Since
    $\E_A(\Id)=\Id$, irreducibility forces the two positive eigenvalues to
    agree, so $|\zeta|=1$ and $Q$ is a positive fixed point of $\E_A$.
    Uniqueness of positive fixed points for an injective normalized tensor then
    makes $Q$ a positive scalar multiple of the identity. Rescaling $X$ by the
    square root of that scalar produces a unitary $V$, and the original
    intertwining relation becomes (\ref{eq:symmetry_phased_covariance}), with
    $|\mu|=1$.
\end{proof}

\begin{theorem}[Virtual unitary gives a twisted-transfer eigenmatrix]%
    \label{thm:twisted_transfer_eigen_of_virtual_unitary}
    \lean{MPSTensor.twistedTransfer_eigen_of_virtualUnitary}
    \leanok
    \uses{def:twisted_transfer_map}
    Assume $\E_A(\Id)=\Id$ and that, for every physical index~$i$,
    $\sum_j u_{ij} A^j = \mu\, V A^i V^\dagger$ with $V$ unitary.
    Then $V$ is an eigenmatrix of the twisted transfer map:
    \begin{align}
        \E_u(V) &= \mu V.
        \label{eq:symmetry_virtual_eigenmatrix}
    \end{align}
\end{theorem}

\begin{proof}
    \leanok
    Expand $\E_u(V)$ using (\ref{eq:symmetry_twisted_transfer}), insert the
    intertwining relation into each local term, and use
    $V^\dagger V = \Id$ together with $\E_A(\Id)=\Id$.
\end{proof}

\begin{theorem}[Boundary-state invariance of the virtual unitary]%
    \label{thm:boundary_state_invariant_of_virtual_unitary}
    \lean{MPSTensor.boundaryState_invariant_of_virtualUnitary}
    \leanok
    \uses{def:injective}
    Let $A$ be injective, assume $u u^\dagger = \Id$, and let $\Lambda$ be a
    positive definite boundary state with $\tr(\Lambda)=1$ and
    $\E_A^\dagger(\Lambda)=\Lambda$.
    If $V$ is a unitary and $\mu$ a phase satisfying, for every physical
    index~$i$,
    \begin{align}
        \sum_j u_{ij} A^j
        &= \mu\, V A^i V^\dagger,
        \label{eq:symmetry_boundary_covariance}
    \end{align}
    then
    \begin{align}
        V^\dagger \Lambda V &= \Lambda.
        \label{eq:symmetry_boundary_invariance}
    \end{align}
\end{theorem}

\begin{proof}
    \leanok
    \uses{thm:psd_fp_unique}
    Relation (\ref{eq:symmetry_boundary_covariance}) shows that
    $V^\dagger \Lambda V$ is another positive fixed point of the adjoint
    transfer map.
    Since $\Lambda$ is the normalized positive fixed point, uniqueness forces
    (\ref{eq:symmetry_boundary_invariance}).
\end{proof}

\begin{lemma}[Local symmetry implies string order]%
    \label{lem:has_string_order_of_local_symmetry}
    \lean{MPSTensor.hasStringOrder_of_localSymmetry}
    \leanok
    \uses{def:local_symmetry, def:has_string_order}
    Assume $\tr(\Lambda)=1$ and $\E_A(\Id)=\Id$.
    If $u$ is a local symmetry of $A$ with respect to $\Lambda$, then string
    order exists for $u$.
\end{lemma}

\begin{proof}
    \leanok
    \uses{thm:twisted_transfer_eigen_of_virtual_unitary,
        def:string_order_boundary_param}
    Choose the boundary matrices $X = V^\dagger$ and $Y = V$, where $V$ is the
    unitary from the local-symmetry datum.
    By Theorem~\ref{thm:twisted_transfer_eigen_of_virtual_unitary},
    one has $\E_u^L(V)=\mu^L V$.
    Therefore
    \begin{align}
        R_L(u;V^\dagger,V)
        &= \mu^L \tr(\Lambda)
         = \mu^L,
        \notag
    \end{align}
    whose norm is constantly $1$.
\end{proof}

\begin{theorem}[Local symmetry iff a unitary peripheral eigenmatrix exists]%
    \label{thm:local_symmetry_iff_spectral_radius}
    \lean{MPSTensor.localSymmetry_iff_spectralRadius_one}
    \leanok
    \uses{def:local_symmetry, def:twisted_transfer_map, def:injective,
        thm:twisted_transfer_spectral_radius_le_one}
    Let $A$ be injective, and assume $u u^\dagger = \Id$, $\E_A(\Id)=\Id$,
    $\Lambda$ is positive definite, $\tr(\Lambda)=1$, and
    $\E_A^\dagger(\Lambda)=\Lambda$.
    Then $u$ is a local symmetry if and only if there exist a unitary
    $V \in \MN{D}$ and a phase $\mu$ with $|\mu|=1$ such that
    $\E_u(V) = \mu V$.
    By Theorem~\ref{thm:twisted_transfer_spectral_radius_le_one}, this witness
    form is equivalent to the peripheral-spectrum condition $\rho(\E_u)=1$.
\end{theorem}

\begin{proof}
    \leanok
    \uses{thm:twisted_transfer_eigen_of_virtual_unitary,
        thm:twisted_transfer_modulus_one_gauge_phase,
        thm:virtual_unitary_of_twisted_gauge_phase,
        thm:boundary_state_invariant_of_virtual_unitary}
    The forward direction is
    Theorem~\ref{thm:twisted_transfer_eigen_of_virtual_unitary}.
    For the reverse direction, start from a unitary peripheral eigenmatrix.
    Theorem~\ref{thm:twisted_transfer_modulus_one_gauge_phase} gives
    gauge-phase equivalence with the companion tensor; then
    Theorem~\ref{thm:virtual_unitary_of_twisted_gauge_phase} yields a virtual
    unitary and phase satisfying the local covariance relation, and
    Theorem~\ref{thm:boundary_state_invariant_of_virtual_unitary} gives the
    boundary-state invariance needed for local symmetry.
\end{proof}

\begin{theorem}[String order $\Leftrightarrow$ local symmetry]%
    \label{thm:string_order_iff_local_symmetry}
    \lean{MPSTensor.stringOrder_iff_localSymmetry}
    \leanok
    \uses{def:has_string_order, def:local_symmetry, def:injective}
    For an injective MPS, assume $u u^\dagger = \Id$,
    $\E_A(\Id)=\Id$, $\Lambda$ is positive definite,
    $\tr(\Lambda)=1$, and $\E_A^\dagger(\Lambda)=\Lambda$. Then string
    order for $u$ exists iff $u$ is a local symmetry.
\end{theorem}

\begin{proof}
    \leanok
    \uses{thm:string_order_implies_twisted_gauge_phase,
        thm:virtual_unitary_of_twisted_gauge_phase,
        thm:boundary_state_invariant_of_virtual_unitary,
        lem:has_string_order_of_local_symmetry}
    If string order exists, Theorem~\ref{thm:string_order_implies_twisted_gauge_phase}
    gives a gauge-phase equivalence with the companion tensor.
    Theorem~\ref{thm:virtual_unitary_of_twisted_gauge_phase} extracts a virtual
    unitary and phase from that equivalence, and
    Theorem~\ref{thm:boundary_state_invariant_of_virtual_unitary} shows that the
    same unitary preserves the boundary state, so one obtains local symmetry.
    Conversely, Lemma~\ref{lem:has_string_order_of_local_symmetry} turns local
    symmetry into a non-decaying boundary string-order sequence.
\end{proof}

\begin{theorem}[Virtual unitary from string order]%
    \label{thm:virtual_unitary_of_string_order}
    \lean{MPSTensor.virtualUnitary_of_stringOrder}
    \leanok
    \uses{def:has_string_order, def:injective}
    For an injective MPS, let $\Lambda \in \MN{D}$, and assume
    $u u^\dagger = \Id$ and $\E_A(\Id)=\Id$.
    If string order exists for $u$ with boundary parameter $\Lambda$, then
    there exists a virtual unitary $V$ and a unit-modulus scalar $\mu$ such that
    $\sum_j u_{ij} A^j = \mu\, V A^i V^\dagger$ for all~$i$.
    Diagrammatically this is the same local intertwining relation,
    but only up to an overall unit scalar on the virtual side:
    % Formula/source: Papers/0802.0447/StringOrder-v10.tex,
    % lines 328--349, gives Condition C1; the theorem's exact formula is
    % sum_j u_{ij}A^j = mu V A^i V^\dagger.
    % Ink-to-index map: i is free above u, j is contracted between u and A,
    % and the two outer virtual indices remain free.
    % Contraction set: one physical bond on the left and two virtual gauge
    % bonds on the right; mu is scalar and adds no index.
    % Type verdict: equality of rank-three local tensors up to a scalar phase.
    % Boundary signature: virtual-west=1 | virtual-east=1 |
    % physical-up=1 | physical-down=0 on both sides.
    \begin{tenkzequation}
        \begin{tenkz}[rows={op:none,ket}]
            \tn[role=operator,up=$i$]{u} \\
            \tn{A}
        \end{tenkz}
        $\;=\;\mu\,$
        \begin{tenkz}[physical=up]
            \tnX{V} & \tn[up=$i$]{A} & \tnX{V^\dagger}
        \end{tenkz}.
    \end{tenkzequation}
\end{theorem}

\begin{proof}
    \leanok
    \uses{thm:string_order_implies_twisted_gauge_phase,
        thm:virtual_unitary_of_twisted_gauge_phase}
    String order first gives gauge-phase equivalence with the companion tensor by
    Theorem~\ref{thm:string_order_implies_twisted_gauge_phase}.
    Theorem~\ref{thm:virtual_unitary_of_twisted_gauge_phase} then normalizes the
    resulting intertwiner to a unitary and produces the phase $\mu$.
\end{proof}

%% ---------------------------------------------------------
